\chapter{実験設定}
% 第5章: 実験
% 本章では実験設定、データセット、比較手法、評価指標について述べる

\section{実験概要}

本章では、第4章で提案したTerrain-Aware A*（TA*）およびField D* Hybridの性能を評価するための実験について述べる。実験の目的は、以下の3点である。第一に、提案手法が既存の経路計画手法と比較して、どの程度の性能向上を実現できるかを定量的に評価する。第二に、地形情報を考慮することによる経路品質の改善効果を検証する。第三に、各手法の計算時間や成功率などの実用性を評価する。

実験環境は以下の通りである。OSはUbuntu 22.04 LTSを使用した。経路計画アルゴリズムはPython 3.10で実装し、数値計算にはNumPy、可視化にはMatplotlibを使用した。計算機は、Intel Core i7プロセッサ（3.6GHz）、メモリ16GBを搭載したワークステーションを使用した。本研究では、アルゴリズムの原理的有効性を検証するため、シミュレーション環境での評価を行った。

\section{ベンチマークデータセット}

\subsection{Dataset3の概要}

実験には、96シナリオから構成されるベンチマークデータセットDataset3を使用した。Dataset3は、経路計画アルゴリズムの性能を多角的に評価するため、マップサイズ（Small 20×20m、Medium 50×50m、Large 100×100m）と地形複雑度（単純<0.15、中程度0.15--0.55、複雑≥0.55）により分類される。各カテゴリ32シナリオ、計96シナリオの統一ベンチマークで評価した。

地形複雑度による分類は、ロボットの走行難易度を段階的に評価するために設計された。単純地形（複雑度<0.15）は平坦で走行しやすく、基本的な経路計画性能を評価する。中程度地形（0.15--0.55）は適度な起伏を含み、地形考慮の効果が現れ始める領域である。複雑地形（≥0.55）は急斜面や粗い表面を含み、地形回避能力が重要となる環境である。

\subsection{Dataset3の生成方法}

Dataset3は、以下の手順でプログラムにより自動生成した。

\textbf{（1）基本地形の生成：}
各カテゴリの特性に応じて、Perlinノイズアルゴリズムを用いて基本地形を生成した。basicカテゴリでは、標高差の小さい平坦な地形を生成した。terrainカテゴリでは、丘、谷、急斜面などの地形変化を含む複雑な地形を生成した。地形の高さは0mから15mの範囲で設定した。

\textbf{（2）障害物の配置：}
obstaclesカテゴリでは、ランダムに障害物を配置した。障害物の数は各シナリオで10個から30個の範囲とし、障害物のサイズは1m×1mから3m×3mの範囲で設定した。障害物は、始点とゴールを完全に遮断しないように配置した。

\textbf{（3）地形複雑度の計算：}
生成された地形に対して、第3章で述べた地形複雑度計算手法を適用した。各セルの標高差、傾斜角、表面粗さを計算し、式(\ref{eq:3_terrain_complexity})により地形複雑度を算出した。

\textbf{（4）始点とゴールの設定：}
各シナリオに対して、始点とゴールを手動で設定した。始点とゴールの距離は15mから25mの範囲とし、直線経路上に地形変化や障害物が存在するように配置した。これにより、単純な直線経路では最適でないシナリオを作成した。

\textbf{（5）妥当性の検証：}
生成されたシナリオに対して、Regular A*を実行し、解が存在することを確認した。解が存在しないシナリオは除外し、再生成を行った。

この生成プロセスにより、再現性のあるベンチマークデータセットを構築した。Dataset3の生成コードは、研究の再現性を確保するため、公開する予定である。

\subsection{マップの詳細設定}

各シナリオのマップは、100×100セルのグリッドで構成され、1セルの解像度は0.2mである。したがって、実空間では20m×20mの領域に相当する。地形情報は、高さ方向に20層のボクセル構造で表現される。各シナリオには、始点とゴールが設定され、アルゴリズムはこれらを結ぶ経路を計画する。

始点とゴールの配置は、各カテゴリの特性に応じて設定される。basicカテゴリでは、始点とゴールが離れており、障害物が少ないため、直線的な経路が期待される。obstaclesカテゴリでは、障害物が始点とゴールの間に配置されており、障害物を迂回する経路が必要となる。terrainカテゴリでは、丘や谷などの地形変化が始点とゴールの間に配置されており、地形を考慮した経路選択が求められる。

\subsection{代表シナリオの紹介}

terrainカテゴリの代表的なシナリオとして、hill\_detour、roughness\_avoidance、combined\_terrainがある。hill\_detourは、始点とゴールの間に丘が配置されており、丘を越えるか迂回するかの判断が求められる。このシナリオは、地形考慮型経路計画の有効性を評価する上で、最も重要なシナリオである。

roughness\_avoidanceは、粗い地形領域が配置されており、表面粗さを考慮した経路選択が重要となる。粗い地形は、ロボットの走行安定性を損なうため、滑らかな地形を選択することが望ましい。combined\_terrainは、複数の地形特性（標高差、傾斜、粗さ）が組み合わさった複雑な環境である。このシナリオでは、すべての地形要素を統合的に評価する能力が試される。

\section{比較手法}

\subsection{比較手法の概要}

本実験では、提案手法の性能を評価するため、4つの経路計画手法を比較した。比較手法は、提案手法2手法（TA*、Field D* Hybrid）と既存手法2手法（AHA*、Theta*）である。公平な比較を行うため、すべての手法に対して同一の環境マップ、始点、ゴールを使用した。また、計算時間の測定条件も統一した。

\subsection{各手法の説明}

比較手法は以下の通りである。

\begin{enumerate}
\item \textbf{TA* (Terrain-Aware A*):} 本研究の提案手法。地形コストを考慮したA*の拡張。式(\ref{eq:4_ta_astar})の評価関数を用いる。地形の走行適性を評価し、安全な経路を生成することを目的とする。

\item \textbf{AHA* (Abstraction Hierarchical A*):} 階層的探索により高速化したA*の拡張手法。環境を階層的に分割し、上位階層で大まかな経路を計画した後、下位階層で詳細化する。計算時間の大幅な短縮が期待できる。

\item \textbf{Theta*:} any-angle経路計画手法。グリッドベースの制約を緩和し、任意の角度での経路生成が可能。Field D*と同様に経路の平滑化が実現されるが、地形情報は考慮しない。

\item \textbf{Field D* Hybrid:} 本研究の提案手法。TA*にField D*の補間技術を統合。地形考慮と経路最適化を同時に実現する。補間により経路を平滑化し、計算効率も向上させる。
\end{enumerate}

\section{評価指標}

\subsection{評価指標の概要}

各手法の性能を評価するため、以下の4つの指標を用いた。これらの指標により、経路品質（経路長、地形コスト）、計算効率（計算時間）、実用性（成功率）を多角的に評価できる。

\subsection{各指標の定義}

\begin{enumerate}
\item \textbf{経路長（Path Length）[m]:} 始点からゴールまでの経路の総距離。短いほど効率的な経路である。経路長は、経路上の各セグメントの距離の総和として計算される。

\item \textbf{計算時間（Computation Time）[s]:} 経路計画にかかった時間。リアルタイム性の指標となる。計算時間は、アルゴリズムの開始から終了までの時間を測定する。

\item \textbf{成功率（Success Rate）[\%]:} 全シナリオのうち、制限時間内に有効な経路を見つけられた割合。実用性の指標となる。成功率は、始点から目標位置への到達が可能であったシナリオ数を全シナリオ数で割った値である。

\item \textbf{地形コスト（Terrain Cost）:} 経路上の地形複雑度の累積値。低いほど走行しやすい経路である。本研究では、点平均法とセグメント加重法の2種類を用いた。
\end{enumerate}

\subsection{地形コストの2つの評価方法}

地形コストの評価には、2つの方法を用いた。第一の方法（点平均法）は、経路上の各ノードの地形コストの平均値である。式(\ref{eq:4_point_average})で定義される。この方法は、経路がどの程度地形の良い領域を通過しているかを直接的に表す。

第二の方法（セグメント加重法）は、各セグメント（ノード間の移動）の距離に応じて地形コストを加重平均する。式(\ref{eq:4_segment_weighted})で定義される。この方法は、実際の移動コストにより近い評価を行える。本研究では点平均法に加えてセグメント加重法を用いた。これは、局所的な凹凸だけでなく、経路全体としての走行負荷を評価するためである。セグメント加重法により、実用的な観点からの経路評価が可能となる。

点平均法は、経路選択のアルゴリズムの性能を直接的に評価するのに適している。一方、セグメント加重法は、実際のロボット走行における総コストを評価するのに適している。両方の方法を用いることで、多角的な評価が可能となる。

\section{実験設定}

\subsection{パラメータ設定}

各手法のパラメータ設定を以下に示す。TA*の地形重み$w_t$は1.0に設定した。これは、第4章で述べた予備実験の結果に基づき、経路長と地形品質のバランスが最良であった値を採用した。ヒューリスティック関数には、ユークリッド距離を用いた。

計算時間の上限は、1シナリオあたり300秒とし、この時間内に経路が見つからない場合は失敗とみなした。また、探索ノード数の上限を100,000ノードに設定し、ノード爆発を防止した。これらの設定により、実用的な範囲での性能評価を行う。

AHA*とTheta*については、デフォルトのパラメータ設定を使用した。Field D* Hybridは、TA*と同じ地形重み$w_t=1.0$を用いた。

\subsection{実験手順}

実験手順は以下の通りである。まず、各手法をDataset3の全96シナリオに対して実行し、経路長、計算時間、成功率、地形コストを記録した。各シナリオは3回実行し、その平均値を採用した。これにより、ランダム性の影響を排除し、安定した評価を行う。

次に、地形複雑度別（単純、中程度、複雑）に各指標の平均値と標準偏差を計算した。これにより、各地形カテゴリにおける手法の性能を比較できる。96シナリオ×3回の実行により、合計288回の経路計画を実施し、各手法の性能を包括的に評価した。

\section{まとめ}

本章では、提案手法の性能を評価するための実験設定について述べた。96シナリオから構成されるDataset3を用いて、4つの経路計画手法を比較した。Dataset3は、Perlinノイズアルゴリズムにより自動生成され、マップサイズと地形複雑度（単純、中程度、複雑）により分類される。評価指標として、経路長、計算時間、成功率、地形コストを用いた。地形コストについては、点平均法とセグメント加重法の2つの評価方法を用いた。次章では、実験結果について詳しく述べる。
